\section{Data}
\label{sec:datasets}

The study uses path and supporting-data versions labeled July~1, 2026. The
path observations cover 00:00--01:00 UTC on that date.

\subsection{RIPE Atlas Measurements}

The corpus contains 18 public RIPE Atlas IPv4 traceroute measurements.
Thirteen cover independently operated anycast DNS Root services; Wikipedia,
Reddit, and \texttt{assets.nflxext.com} provide three application-facing
measurements; measurements 5051 and 5151 use dynamic targets as multi-target
UDP and ICMP references. Table~\ref{tab:atlas_measurements} reports the corpus
after canonical trace-identity deduplication.

\begin{table*}[t]
  \centering
  \caption{RIPE Atlas measurement corpus.}
  \label{tab:atlas_measurements}
  \footnotesize
  \setlength{\tabcolsep}{4.2pt}
  \begin{tabular}{p{0.15\textwidth}p{0.27\textwidth}p{0.30\textwidth}rr}
    \toprule
    Group & Target & MSM ID(s) & Raw & Valid \\
    \midrule
    DNS Roots & A--M Root
      & 5009, 5010, 5011, 5012, 5013, 5004, 5014, 5015, 5005, 5016,
        5001, 5008, 5006
      & 534,276 & 317,873 \\
    Applications & Wikipedia; Reddit; Netflix Assets
      & 86710103; 176906957; 176517335
      & 115,596 & 115,596 \\
    Topology references & IPv4 UDP; IPv4 ICMP
      & 5051; 5151
      & 57,442 & 57,442 \\
    \midrule
    \textbf{Total} & \textbf{18 measurements} & --
      & \textbf{707,314} & \textbf{490,911} \\
    \bottomrule
  \end{tabular}
\end{table*}

Each record retains the destination address returned by RIPE Atlas. This
preserves both the multi-target structure of the two dynamic-target
measurements and the destination instances selected by the service
measurements during the window.

The three families are not a controlled comparison. The dynamic-target
measurements sample many destinations, whereas each service measurement
samples the destinations selected for that service. We use the former as
topology references and do not attribute family differences to service
deployment or target selection.

\subsection{Supporting Data}

RIPE Atlas probe metadata supplies probe country and source ASN
~\cite{ripe_atlas}. IPinfo Location and ASN databases annotate visible hops
~\cite{ipinfo}. The Submarine Cable Map provides 694 cables and 1,916 landing
points, including lifecycle and organizational fields
~\cite{telegeography_submarine_map}. CAIDA AS Relationships supplies
supplementary AS context~\cite{caida_as_relationships}. A versioned country
taxonomy classifies probe countries as island or archipelagic, landlocked, or
other coastal.

\begin{table*}[t]
  \centering
  \caption{Supporting datasets.}
  \label{tab:supporting_datasets}
  \footnotesize
  \begin{tabular}{p{0.30\textwidth}p{0.25\textwidth}p{0.35\textwidth}}
    \toprule
    Dataset & Version label & Content \\
    \midrule
    RIPE Atlas probe metadata & 2026-07-01 & Probe country and source ASN \\
    IPinfo Location MMDB & 2026-07-01 & Hop geolocation \\
    IPinfo ASN MMDB & 2026-07-01 & IP--AS mapping \\
    Submarine Cable Map & 2026-07-01 & 694 cables; 1,916 landing points \\
    CAIDA AS Relationships & 2026-07-01 & AS relationships \\
    Country geography taxonomy & 2026-07-01 & Operational geography classes \\
    \bottomrule
  \end{tabular}
\end{table*}
